\section*{Exercises about Dimension}

\exercise{1}
\begin{xrcs}
  Suppose $V$ is finite-dimensional and $U$ is a subspace of $V$ such that $\dim U = \dim V$. Prove that $U = V$.

  \begin{xprf}
    Let $u_1, \dots, u_n$ be a basis of $U$, so $n = \dim U = \dim V$.
    The list $u_1, \dots, u_n$ is linearly independent in $V$ and has length $n = \dim V$.
    By Theorem \ref{thm: linearly independent list of the right length is a basis}, every linearly independent list of length $\dim V$ in a finite-dimensional space $V$ is a basis of $V$.
    Therefore, $V = \myspan{u_1, \dots, u_n} = U$.
  \end{xprf}
\end{xrcs}

\begin{ainote}[The finite-dimensional hypothesis is doing real work]
  For infinite-dimensional $V$ the statement fails, and it fails badly. Take $V = \polyn(\real)$, the space of all polynomials, and let $U$ be the subspace of polynomials with zero constant term. Both are infinite-dimensional, so \qt{$\dim U = \dim V$} holds in the only sense available, yet $U \neq V$ since $1 \notin U$. The proof above breaks at the very first move: $U$ has no finite basis $u_1, \ddd, u_n$ to extract.

  The finite-dimensional analogue of that failure is worth keeping in mind too: a proper subspace always has \emph{strictly} smaller dimension, so $\dim U = \dim V$ combined with $U \subseteq V$ is a genuinely strong hypothesis, not a formality.
\end{ainote}

\exercise{8}
\begin{xrcs}
  Suppose $U$ is a subspace of a finite-dimensional vector space $V$. Prove that there exists a subspace $W$ of $V$ such that
  \begin{equation}
    V = U \oplus W.
  \end{equation}
  \begin{xprf}
    Let $u_1, \dots, u_m$ be a basis of $U$.
    By Theorem \ref{thm: every linearly independent list of vectors in a finite-dimensional vector space can be extended to a basis of the vector space}, this list can be extended to a basis
    \[
      u_1, \dots, u_m, w_1, \dots, w_k
    \]
    of $V$, where $m + k = \dim V$.
    Define $W := \myspan{w_1, \dots, w_k}$.
    Since the combined list spans $V$, $V = U + W$.
    Moreover, if $v \in U \cap W$, then $v = a_1 u_1 + \dots + a_m u_m = b_1 w_1 + \dots + b_k w_k$, so
    $a_1 u_1 + \dots + a_m u_m - b_1 w_1 - \dots - b_k w_k = 0$.
    By linear independence of the basis, all coefficients are $0$, hence $v = 0$.
    Therefore, $U \cap W = \{0\}$, which proves $V = U \oplus W$.
  \end{xprf}
\end{xrcs}

\begin{ainote}[The complement $W$ is never unique]
  The exercise asks for \emph{a} subspace $W$, and that indefinite article is deliberate: as soon as $U \neq \{0\}$ and $U \neq V$, there are many such $W$, one for essentially every way of choosing the extension vectors $w_1, \ddd, w_k$. Exercise 1C.23 is exactly this phenomenon written out in $\real^2$, where the $x$-axis has both the $y$-axis and the line $y = x$ as complements.

  So $W$ is not \qt{the complement of $U$} in any canonical sense, and the notation $V = U \oplus W$ should not be read as defining $W$ from $U$. An inner product is what later supplies a canonical choice, namely $W = U^{\perp}$; without one, no choice is distinguished.
\end{ainote}

\exercise{19}
\begin{xrcs}
  Explain why you might guess, motivated by analogy with the formula for the number of elements in the union of three finite sets, that if $V_1$, $V_2$, $V_3$ are subspaces of a finite-dimensional vector space $V$, then
  \begin{equation}
      \begin{aligned}
      \dim (V_1 + V_2 + V_3) &= \dim (V_1) + \dim (V_2) + \dim (V_3) \\
      & \quad - \dim(V_1 \cap V_2) - \dim (V_1 \cap V_3) - \dim (V_2 \cap V_3) \\
      & \quad + \dim (V_1 \cap V_2 \cap V_3).
    \end{aligned}
  \end{equation}


  Then either prove the formula above or give a counterexample.

  \begin{xsol}
    The formula is \textbf{wrong}. Consider the following subspaces of $\real^2$:
    \begin{itemize}
      \item $V_1 := \{ (x,0) \mid x \in \real \} $ \quad (the $x$-axis)
      \item $V_2 := \{ (x,x) \mid x \in \real \} $ \quad (the line $y=x$)
      \item $V_3 := \{ (0,y) \mid y \in \real \} $ \quad (the $y$-axis)
    \end{itemize}

    The right-hand side of the formula gives: $1+1+1-0-0-0+0 = 3$, whereas the left-hand side is $\dim (V_1 + V_2 + V_3) = 2$. So the formula \emph{does not hold} for this example.

    However, \emph{it does hold}  in every case where one of the subspaces is contained in another; hence, without loss of generality, assume that $V_1 \subseteq V_3$. For this situation, we have the following lemma.

    \prooffont{Lemma:} If $V_1, V_2, V_3$ are finite-dimensional subspaces of a vector space, and $V_1 \subseteq V_3$, then
    \begin{equation}
      \label{eq: lemma}
       (V_1 + V_2) \cap V_3 = (V_1 \cap V_3) + (V_2 \cap V_3)
    \end{equation}

    \begin{xprf}
      If $v \in (V_1 \cap V_3) + (V_2 \cap V_3)$, then we can write $v = u + w$, where $u \in V_1\cap V_3$ and  $w \in V_2 \cap V_3$.
      Because $V_3$ is closed under addition and both $u$ and $w$ are in $V_3$, so is $v$. Hence, $v \in V_3$ and $v=u+w$, such that $u \in V_1$ and $w \in V_2$. Therefore, $v \in (V_1 + V_2) \cap V_3$. We conclude:
      \begin{equation}
        \label{eq: first inclusion}
        (V_1 \cap V_3) + (V_2 \cap V_3) \subseteq (V_1 + V_2) \cap V_3
      \end{equation}

      Conversely, if $v \in (V_1 + V_2) \cap V_3$, then we can write $v = u + w$, where $u \in V_1$, $w \in V_2$, and $v \in V_3$. Since $V_1 \subseteq V_3$ (by our assumption), we can rewrite
      \begin{equation}
        w = \underbrace{v}_{\in \; V_3}-\underbrace{u}_{\in \; V_3}.
      \end{equation}

      Since both $u$ and $v$ belong to $V_3$, so does $w$. Therefore, $v=u+w$ such that $u \in V_1 \cap V_3$ and $w \in V_2 \cap V_3$. Hence, $v \in (V_1 \cap V_3) + (V_2 \cap V_3)$. We conclude:
      \begin{equation}
        \label{eq: second inclusion}
        (V_1 + V_2) \cap V_3 \subseteq (V_1 \cap V_3) + (V_2 \cap V_3)
      \end{equation}

      Inclusion \eqref{eq: first inclusion} and \eqref{eq: second inclusion} together yield
      \begin{equation}
        (V_1 + V_2) \cap V_3 = (V_1 \cap V_3) + (V_2 \cap V_3)
      \end{equation}

      for $V_1 \subseteq V_3$.
    \end{xprf}

    Using the standard dimension formula  for the sum of two subspaces \ref{thm: dimension of a sum of subspaces}, we obtain
    \begin{equation}
      \label{eq: first forumla}
      \begin{aligned}
        \dim (V_1 + V_2 + V_3) &= \dim ((V_1+V_2) + V_3) \\
          &= \dim (V_1 + V_2) + \dim (V_3) - \dim((V_1 +V_2) \cap V_3) \\
          &= \dim (V_1) + \dim (V_2) + \dim (V_3) \\
          & \quad - \dim (V_1 \cap V_2) - \dim ((V_1 + V_2) \cap V_3)
      \end{aligned}
    \end{equation}

    Using the assumption that $V_1 \subseteq V_3$, we calculate:
    \begin{equation}
      \begin{aligned}
        \dim ((V_1 + V_2) \cap V_3)
        &= \dim ((V_1 \cap V_3) + (V_2 \cap V_3)) \\
        &= \dim (V_1 \cap V_3) + \dim (V_2 \cap V_3) - \underbrace{\dim ((V_1 \cap V_3) \cap (V_2\cap V_3))}_{= \; \dim ((V_1 \cap V_2 \cap V_3))}
      \end{aligned}
    \end{equation}

    If we plug in this result in the first formula \eqref{eq: first forumla}, we have
    \begin{equation}
      \begin{aligned}
        \dim (V_1 + V_2 + V_3)
        &= \dim (V_1) + \dim (V_2) + \dim (V_3) - \dim (V_1 \cap V_2) - \dim ((V_1 + V_2) \cap V_3) \\
        &= \dim (V_1) + \dim (V_2) + \dim (V_3) - \dim (V_1 \cap V_2) \\
        &\quad - \big(\dim (V_1 \cap V_3) + \dim (V_2 \cap V_3) - \dim ((V_1 \cap V_2 \cap V_3)) \big) \\
        &= \dim (V_1) + \dim (V_2) + \dim (V_3) \\
        & \quad - \dim(V_1 \cap V_2) - \dim (V_1 \cap V_3) - \dim (V_2 \cap V_3) \\
        & \quad + \dim (V_1 \cap V_2 \cap V_3).
      \end{aligned}
    \end{equation}

    So, under certain conditions (i.e. $V_1 \subseteq V_3$), the formula is \textbf{true}.
  \end{xsol}
\end{xrcs}

\begin{xrcs}
  Show that for a finite-dimensional vector space $V$ with subspaces $V_1$, $V_2$, $V_3$ the following formula holds:
  \begin{equation}
    \begin{aligned}
      \dim (V_1 + V_2 + V_3)  &= \dim (V_1) + \dim (V_2) + \dim (V_3) \\
      & \qquad - \frac{\dim (V_1 \cap V_2) + \dim (V_1 \cap V_3) + \dim (V_2 \cap V_3)}{3} \\
      & \qquad - \frac{ \dim \left(  (V_1 + V_2) \cap V_3 \right) + \dim \left( (V_1 + V_3) \cap V_2 \right)+ \dim \left( (V_2 + V_3) \cap V_1 \right) }{3} \\
    \end{aligned}
  \end{equation}

  \prooffont{Solution:} Recall that for any two subspaces $U$ and $W$ of a finite-dimensional vector space we have
  \begin{equation}
    \dim (U + W)  = \dim U + \dim W - \dim (U \cap W).
  \end{equation}

  So we can use associativity and commutativity of the sum $V_1 + V_2 + V_3$. In particular:
  \begin{equation}
    \begin{aligned}
      V_1 + V_2 + V_3 & = (V_1 + V_2) + V_3 \\
                      & = V_1 + (V_2 + V_3) \\
                      & = (V_1 + V_3) + V_2
    \end{aligned}
  \end{equation}

  \begin{equation}
    \begin{aligned}
      \dim((V_1 + V_2) + V_3)
      & = \dim(V_1 + V_2) + \dim(V_3) - \dim ((V_1 + V_2) \cap V_3) \\
      & = \dim V_1 + \dim V_2 + \dim V_3 - \dim(V_1 \cap V_2) - \dim ((V_1 + V_2) \cap V_3) \\
      \dim(V_1 + (V_2 + V_3))
      & = \dim(V_2 + V_3) + \dim(V_1) - \dim ((V_2 + V_3) \cap V_1) \\
      & = \dim V_1 + \dim V_2 + \dim V_3 - \dim(V_2 \cap V_3) - \dim ((V_2 + V_3) \cap V_1) \\
      \dim((V_1 + V_3) + V_2)
      & = \dim(V_1 + V_3) + \dim(V_2) - \dim ((V_1 + V_3) \cap V_2) \\
      & = \dim V_1 + \dim V_2 + \dim V_3 - \dim(V_1 \cap V_3) - \dim ((V_1 + V_3) \cap V_2) \\
    \end{aligned}
  \end{equation}

  Therefore, adding the three identities yields
  \begin{equation}
    \begin{aligned}
      3 \dim (V_1 + V_2 + V_3) &=
      \dim V_1 + \dim V_2 + \dim V_3 - \dim(V_1 \cap V_2) - \dim ((V_1 + V_2) \cap V_3) \\
      & \quad + \dim V_1 + \dim V_2 + \dim V_3 - \dim(V_2 \cap V_3) - \dim ((V_2 + V_3) \cap V_1) \\
      & \quad + \dim V_1 + \dim V_2 + \dim V_3 - \dim(V_1 \cap V_3) - \dim ((V_1 + V_3) \cap V_2) \\
      & = 3 \dim V_1 + 3 \dim V_2 + 3 \dim V_3 \\
      & \quad - \dim(V_1 \cap V_2) - \dim(V_2 \cap V_3) - \dim(V_1 \cap V_3) \\
      & \quad - \dim ((V_1 + V_2) \cap V_3) - \dim ((V_2 + V_3) \cap V_1) - \dim ((V_1 + V_3) \cap V_2)
    \end{aligned}
  \end{equation}

  Dividing both sides by $3$ gives the formula in question.
\end{xrcs}

\begin{ainote}[Why this averaging trick works, and what it does not give you]
  The three displayed identities are the \emph{same} theorem \ref{thm: dimension of a sum of subspaces} applied three times, once for each way of bracketing $V_1 + V_2 + V_3$ into a sum of two subspaces. Each one is correct on its own, so each already computes $\dim(V_1+V_2+V_3)$ exactly. Adding them and dividing by $3$ therefore proves a true formula, but it buys no new information: it is the average of three correct statements, deliberately made symmetric in $V_1, V_2, V_3$.

  That symmetry is the entire point of the exercise. The naive inclusion--exclusion guess of Exercise 19 is symmetric but \emph{false}; the bracketed formula \eqref{eq: first forumla} is true but breaks the symmetry by singling out one subspace. Averaging over all three bracketings restores the symmetry without giving up correctness.

  What the formula still does not do is express $\dim(V_1+V_2+V_3)$ purely in terms of the dimensions of the $V_k$ and of their pairwise and triple intersections. The mixed terms $\dim((V_1+V_2) \cap V_3)$ survive on the right-hand side, and by Exercise 19 they genuinely cannot be reduced to intersection dimensions in general.
\end{ainote}

\begin{failedattempt}[Alternative attempt via basis extension (does not work)]
  Let $r_1, \ddd, r_{n_{123}} \in V$ be a basis of $V_1 \cap V_2 \cap V_3$, where $n_{123} \in \nat_{\geq 0}$, meaning that the basis can also be the empty list and $V_1 \cap V_2 \cap V_3$ might be equal to $\{\vec 0\}$. Note that the empty list is defined to be linearly independent and that its span is defined to be $\{\vec 0\}$. By this definition, we have
  \begin{equation}
    \label{eq: dim for V_1 cap V_2 cap V_3}
    \dim (V_1 \cap V_2 \cap V_3) = n_{123}.
  \end{equation}

  Since $r_1, \ddd, r_{n_{123}}$ is a linearly independent list in $V$ (as well as in $V_1 \cap V_2$, $V_1 \cap V_3$, and $V_2 \cap V_3$), this list can be extended to a basis of $V_1 \cap V_2$, a basis of $V_1 \cap V_2$, and a basis $V_2 \cap V_3$ respectively. So let
  \begin{itemize}
    \item $r_1, \ddd, r_{n_{123}}, v_1, \ddd, v_{m_{12}}$ be a basis of $V_1 \cap V_2$,
    \item $r_1, \ddd, r_{n_{123}}, \widetilde{v}_1, \ddd, \widetilde{v}_{m_{13}}$ be a basis of $V_1 \cap V_3$ and
    \item $r_1, \ddd, r_{n_{123}}, \overline{v}_1, \ddd, \overline{v}_{m_{23}}$ be a basis of $V_2 \cap V_3$ for some $m_{12}, m_{13}, m_{23} \in \nat_{\geq 0}$.
  \end{itemize}

  Thus, we have
  \begin{equation}
    \label{eq: dim for V_1 cap V_2, V_1 cap V_3, V_2 cap V_3}
    \dim (V_1 \cap V_2) = n_{123} + m_{12}, \; \dim (V_1 \cap V_3) = n_{123} + m_{13}, \myand \dim (V_2 \cap V_3) = n_{123} + m_{23}.
  \end{equation}

  Note that the extensions might be empty, and $m_{12}, m_{13}, m_{23} \in \nat_{\geq 0}$, meaning $V_1 \cap V_2$, $V_1 \cap V_3$ and $V_2 \cap V_3$ could be $V_1 \cap V_2 \cap V_3$. This further implies that
  \begin{itemize}
    \item $v_1, \ddd, v_{m_{12}}$ is a basis of $(V_1 \cap V_2) \backslash (V_1 \cap V_2 \cap V_3)$,
    \item $\widetilde{v}_1, \ddd, \widetilde{v}_{m_{13}}$ is a basis of $(V_1 \cap V_3) \backslash (V_1 \cap V_2 \cap V_3)$, and
    \item $\overline{v}_1, \ddd, \overline{v}_{m_{23}}$ is a basis of $(V_2 \cap V_3) \backslash (V_1 \cap V_2 \cap V_3)$.
  \end{itemize}
  Let us do a simple check. Suppose $V_1 \cap V_2 =  V_1 \cap V_2 \cap V_3$. This would mean that
  $(V_1 \cap V_2) \backslash (V_1 \cap V_2 \cap V_3) = \{ \vec 0 \}$, so $m_{12} = 0$. The empty list is a basis of $\{ \vec 0 \}$, so everything is fine. Note that there is a small difference compared to set notation here. If we write $U \backslash W$ for two vector spaces $U$ and $W$, we mean the subspace of $U$ that is linearly independent of $W$. So if we have $W \backslash W$, the only subspace of $W$ that is linearly independent of all vectors in $W$ is $\{ \vec 0 \}$. In set notation we would have $W \backslash W = \varnothing$.

  Later on we will use the fact that
  \begin{equation}
    \label{eq: fact for dim}
    \dim (V_1 \cap V_2) + \dim (V_1 \cap V_3) + \dim (V_2 \cap V_3) = 3\cdot n_{123} + 2\cdot (m_{12} + m_{13} + m_{23}).
  \end{equation}

  Basis of $V_1 + V_2$ looks like this:
  \[
    r_1, \ddd, r_{n_{123}}, v_1, \ddd, v_{m_{12}}, \widetilde{v}_1, \ddd, \widetilde{v}_{m_{13}}, u_1, \ddd, u_j, \overline{v}_1, \ddd, \overline{v}_{m_{23}}, w_1, \ddd, w_{m_{23}}
  \]

  So let us further extend
  \begin{itemize}
    \item $r_1, \ddd, r_{n_{123}}, v_1, \ddd, v_{m_{12}}, \widetilde{v}_1, \ddd, \widetilde{v}_{m_{13}}, u_1, \ddd, u_j$ to a basis of $V_1$,
    \item $r_1, \ddd, r_{n_{123}}, v_1, \ddd, v_{m_{12}}, \overline{v}_1, \ddd, \overline{v}_{m_{23}}, w_1, \ddd, w_k$ to a basis of $V_2$ and
    \item $r_1, \ddd, r_{n_{123}}, \widetilde{v}_1, \ddd, \widetilde{v}_{m_{13}}, \overline{v}_1, \ddd, \overline{v}_{m_{23}}, z_1, \ddd, z_l$ to a basis of $V_3$ for $j,k,l \in \nat_{\geq 0}$.
  \end{itemize}

  Again there could be no $u$'s, $w$'s or $z$'s to be added. This will yield
  \begin{itemize}
    \item $\dim (V_1) = n_{123} + m_{12} + m_{13} + j$,
    \item $\dim (V_2) = n_{123} + m_{12} + m_{23} + k$, and
    \item $\dim (V_3) = n_{123} + m_{13} + m_{23} + l$.
  \end{itemize}

  Adding these equations together we will get
  \begin{equation}
    \label{eq: dim V_1 + dim V_2 + dim V_3}
      \dim (V_1)+ \dim (V_2) + \dim (V_3) = \Big(3 \cdot n_{123} + 2 \cdot (m_{12}+m_{13}+m_{23})+j+k+l \Big)
  \end{equation}

  Before we can proceed we need to justify that
  \begin{itemize}
    \item $r_1,  \ddd r_{n_{123}} ,  \widetilde{v}_1, \ddd, \widetilde{v}_{m_{13}}, \overline{v}_1, \ddd, \overline{v}_{m_{23}}$ is a basis of $(V_1+V_2) \cap V_3$, hence $\dim ((V_1+V_2) \cap V_3) = n_{123} + m_{13}+m_{23}$.
    \item $r_1, \ddd, r_{n_{123}}, v_1, \ddd, v_{m_{12}}, \overline{v}_1, \ddd, \overline{v}_{m_{23}}$ is a basis of $(V_1+V_3) \cap V_2$, hence $\dim ((V_1+V_3) \cap V_2) = n_{123} + m_{12}+m_{23}$.
    \item $r_1, \ddd, r_{n_{123}}, v_1, \ddd, v_{m_{12}}, \widetilde{v}_1, \ddd, \widetilde{v}_{m_{13}}$ is a basis of $(V_2+V_3) \cap V_1$, hence $\dim ((V_2+V_3) \cap V_1) = n_{123} + m_{12}+m_{13}$.
  \end{itemize}

  So we need to show that $r_1,  \ddd r_{n_{123}} ,  \widetilde{v}_1, \ddd, \widetilde{v}_{m_{13}}, \overline{v}_1, \ddd, \overline{v}_{m_{23}}$ is a basis of $(V_1+V_2) \cap V_3$.
  Looking at the basis of $V_3$, for $u \in (V_1 + V_2) \cap V_3$ we can write $u$ as
  \begin{equation}
    u = (\alpha'_1 r_1+ \cdots+ \alpha'_{n_{123}} r_{n_{123}})
    + (\beta'_1 \widetilde{v}_1+ \cdots+ \beta'_{m_{13}} \widetilde{v}_{m_{13}})
    + (\gamma'_1 \overline{v}_1+ \cdots+ \gamma'_{m_{23}} \overline{v}_{m_{23}})
    + (\delta'_1 z_1+\cdots+ \delta'_l z_l)
  \end{equation}

  If we can show that
  \begin{equation}
    \label{eq: basis of V_1 + V_2 + V_3}
    r_1, \ddd, r_{n_{123}}, v_1, \ddd, v_{m_{12}}, \widetilde{v}_1, \ddd, \widetilde{v}_{m_{13}}, \overline{v}_1, \ddd, \overline{v}_{m_{23}}, u_1, \ddd, u_j, w_1, \ddd, w_k, z_1, \ddd, z_l
  \end{equation}
  is a basis of $V_1 + V_2 + V_3$, then
  \begin{equation}
    \dim (V_1 + V_2 + V_3) = n_{123} + m_{12} + m_{13} + m_{23} + j + k + l
  \end{equation}

  Using \eqref{eq: dim V_1 + dim V_2 + dim V_3}, \eqref{eq: fact for dim}, and \eqref{eq: dim for V_1 cap V_2, V_1 cap V_3, V_2 cap V_3}:
  \[
  \begin{aligned}
    \dim (V_1 + V_2 + V_3) & = n_{123} + m_{12} + m_{13} + m_{23} + j + k + l \\
    & = \Big(3 \cdot n_{123} + 2 \cdot (m_{12}+m_{13}+m_{23})+j+k+l \Big) \\
    & \qquad - \frac{3 \cdot n_{123}+(m_{12}+m_{13}+m_{23}) }{3} \\
    & \qquad - \frac{(n_{123} + m_{13} + m_{23})+(n_{123} + m_{12} + m_{23})+(n_{123} + m_{12} + m_{13})}{3} \\
    & = \dim (V_1) + \dim (V_2) + \dim (V_3) \\
    & \qquad - \frac{\dim (V_1 \cap V_2) + \dim (V_1 \cap V_3) + \dim (V_2 \cap V_3)}{3} \\
    & \qquad - \frac{ \dim \left(  (V_1 + V_2) \cap V_3 \right) + \dim \left( (V_1 + V_3) \cap V_2 \right)+ \dim \left( (V_2 + V_3) \cap V_1 \right) }{3}.
  \end{aligned}
  \]

  Now let all $a$'s, $b$'s, $c$'s, $d$'s, $e$'s, $f$'s and $g$'s be elements of $\myF$ such that
  \begin{equation}
    \label{eq: linear independence formula}
    \begin{aligned}
    0 = \underbrace{a_1r_1 + \cdots + a_{n_{123}} r_{n_{123}}}_{\substack{\in \; V_1 \cap V_2 \cap V_3 \\ (\text{common intersection})}}
    + \underbrace{b_1v_1 + \cdots + b_{m_{12}} v_{m_{12}}}_{
        \substack{\in \; (V_1 \cap V_2) \backslash (V_1 \cap V_2 \cap V_3) \\ (\text{extension for } V_1 \cap V_2)}}
    + \underbrace{c_1\widetilde{v}_1 + \cdots + c_{m_{13}}\widetilde{v}_{m_{13}}}_{
        \substack{\in \; (V_1 \cap V_3) \backslash (V_1 \cap V_2 \cap V_3) \\ (\text{extension for } V_1 \cap V_3)}}
    + \underbrace{d_1\overline{v}_1 + \cdots + d_{m_{23}}\overline{v}_{m_{23}}}_{
        \substack{\in \; (V_2 \cap V_3) \backslash (V_1 \cap V_2 \cap V_3) \\ (\text{extension for } V_2 \cap V_3)}} \\
    \quad
    + \underbrace{e_1 u_1 + \cdots + e_j u_j}_{
      \substack{\in \; V_1 \\ (\text{extension for } V_1)}}
    + \underbrace{f_1 w_1 + \cdots + f_k w_k \phantom{_j}\!\!\!}_{
      \substack{\in \; V_2 \\ (\text{extension for } V_2)}}
    + \underbrace{g_1 z_1 + \cdots + g_l z_l \phantom{_j}\!\!\!}_{
      \substack{\in \; V_3 \\ (\text{extension for } V_3)}}
    \end{aligned}
  \end{equation}

  Isolating the $u$ terms:
  \[
  \begin{aligned}
    e_{1} u_1 + \cdots + e_j u_j = &- a_1r_1 -  \cdots - a_{n_{123}} r_{n_{123}} - b_1v_1 -  \cdots - b_{m_{12}} v_{m_{12}} \\
    &-  c_1\widetilde{v}_1 - \cdots - c_{m_{13}}\widetilde{v}_{m_{13}} - d_{1}\overline{v}_1 - \cdots - d_{m_{23}}\overline{v}_{m_{23}} \\
    &  - f_1 w_1 - \cdots - f_k w_k - g_1 z_1 - \cdots - g_l z_l
  \end{aligned}
  \]
  which shows that $e_{1} u_1 + \cdots + e_j u_j \in V_1 \cap (V_2 + V_3)$.
  Because $r_{1}, \ddd, r_{n_{123}}, v_{1}, \ddd, v_{m_{12}}, \widetilde{v}_1, \ddd, \widetilde{v}_{m_{13}}$ is a basis of $V_1 \cap (V_2+V_3)$,
  linear independence of the basis of $V_1$ forces all $e$'s to be $0$.
  Analogously, all $f$'s and $g$'s are $0$.
  Then \eqref{eq: linear independence formula} simplifies to:
  \begin{equation}
    \label{eq: second linear independence formula}
    a_1r_1 +  \cdots + a_{n_{123}} r_{n_{123}} + b_1v_1 +  \cdots + b_{m_{12}} v_{m_{12}} +  c_1\widetilde{v}_1 + \cdots + c_{m_{13}}\widetilde{v}_{m_{13}} + d_{1}\overline{v}_1 + \cdots + d_{m_{23}}\overline{v}_{m_{23}} = 0.
  \end{equation}
  Isolating $b_1 v_1 + \dots + b_{m_{12}} v_{m_{12}}$ and using the basis properties shows that all $b$'s, $a$'s, $c$'s, and $d$'s are $0$.
  Thus, the list is linearly independent and forms a basis of $V_1 + V_2 + V_3$.
\end{failedattempt}

\begin{ainote}[Where it breaks]
  The construction is genuinely clever, and everything up to the dimension bookkeeping for $V_1, V_2, V_3$ is fine. The break is the innocuous-looking sentence \qt{Before we can proceed we need to justify that $r_1, \ddd, r_{n_{123}}, \widetilde{v}_1, \ddd, \widetilde{v}_{m_{13}}, \overline{v}_1, \ddd, \overline{v}_{m_{23}}$ is a basis of $(V_1+V_2) \cap V_3$.}

  That claim says the part of $V_3$ reachable from $V_1 + V_2$ is spanned by the vectors coming from $V_1 \cap V_3$ and from $V_2 \cap V_3$. Written without bases, it is exactly
  \begin{equation}
    (V_1+V_2) \cap V_3 = (V_1 \cap V_3) + (V_2 \cap V_3),
  \end{equation}
  the \emph{modular identity} --- and \eqref{eq: lemma} above proves it only under the extra hypothesis $V_1 \subseteq V_3$. Without such a containment it is false, as the very counterexample opening Exercise 19 shows: with $V_1$ the $x$-axis, $V_2$ the line $y=x$ and $V_3$ the $y$-axis in $\real^2$,
  \begin{equation}
    (V_1 + V_2) \cap V_3 = \real^2 \cap V_3 = V_3, \quad \text{but} \quad (V_1 \cap V_3) + (V_2 \cap V_3) = \{0\} + \{0\} = \{0\}.
  \end{equation}

  So a vector of $(V_1+V_2) \cap V_3$ need not decompose into a piece from $V_1 \cap V_3$ plus a piece from $V_2 \cap V_3$, and the proposed list simply fails to span $(V_1+V_2) \cap V_3$. Since the whole computation rests on that dimension count, the construction cannot prove the general formula. This is not a repairable slip: by Exercise 19 no symmetric formula in the intersection dimensions alone can be correct, so any basis-extension argument organised purely around $V_1 \cap V_2 \cap V_3$ and the pairwise intersections is doomed before it starts.

  Kept here because the failure is instructive: bases built by repeated extension \emph{look} like they compose, and the assumption that they do is one of the easiest errors to make invisibly.
\end{ainote}

\begin{failedattempt}[Attempt to prove modularity for \emph{all} subspaces (does not work)]
  \prooffont{Statement.} We are asked to prove the following equality between subspaces of a vector space:
  \[
    (V_1 + V_2) \cap V_3 = (V_1 \cap V_3) + (V_2 \cap V_3),
  \]
  where $V_1, V_2, V_3$ are subspaces of some vector space, with \emph{no} containment assumption.

  \prooffont{Inclusion 1: $(V_1 + V_2) \cap V_3 \subseteq (V_1 \cap V_3) + (V_2 \cap V_3)$.}
  Let $v \in (V_1 + V_2) \cap V_3$. This means that $v$ satisfies two conditions: $v \in V_1 + V_2$, which implies that $v$ can be written as $v = v_1 + v_2$ for some $v_1 \in V_1$ and $v_2 \in V_2$; and $v \in V_3$. Thus
  \[
    v = v_1 + v_2 \quad \text{with} \quad v_1 \in V_1, \, v_2 \in V_2, \, v \in V_3.
  \]
  We aim to show that $v_1 \in V_1 \cap V_3$ and $v_2 \in V_2 \cap V_3$, which would give $v \in (V_1 \cap V_3) + (V_2 \cap V_3)$.

  Now, suppose for the sake of contradiction that $v_1 \notin V_3$ or $v_2 \notin V_3$. Since $v_1 \in V_1$, $v_2 \in V_2$, and $v = v_1 + v_2 \in V_3$, we note that $V_3$ is a subspace and therefore closed under addition. If either $v_1 \notin V_3$ or $v_2 \notin V_3$, their sum cannot be in $V_3$, leading to a contradiction. Therefore, both $v_1 \in V_3$ and $v_2 \in V_3$ must hold.

  Thus $v_1 \in V_1 \cap V_3$ and $v_2 \in V_2 \cap V_3$, and consequently $v \in (V_1 \cap V_3) + (V_2 \cap V_3)$.

  \prooffont{Inclusion 2: $(V_1 \cap V_3) + (V_2 \cap V_3) \subseteq (V_1 + V_2) \cap V_3$.}
  Let $v \in (V_1 \cap V_3) + (V_2 \cap V_3)$, so $v = v_1 + v_2$ with $v_1 \in V_1 \cap V_3$ and $v_2 \in V_2 \cap V_3$. Then $v_1 \in V_1$ and $v_2 \in V_2$ give $v \in V_1 + V_2$, while $v_1, v_2 \in V_3$ and closedness of $V_3$ under addition give $v \in V_3$. Hence $v \in (V_1 + V_2) \cap V_3$.

  \prooffont{Conclusion.} Both inclusions hold, therefore
  \[
    (V_1 + V_2) \cap V_3 = (V_1 \cap V_3) + (V_2 \cap V_3),
  \]
  which completes the proof.
\end{failedattempt}

\begin{ainote}[Where it breaks]
  Inclusion 2 is correct and needs no hypothesis at all; it is the same argument as in \eqref{eq: first inclusion}. Everything hangs on Inclusion 1, and there the sentence
  \begin{center}
    \qt{If either $v_1 \notin V_3$ or $v_2 \notin V_3$, their sum cannot be in $V_3$}
  \end{center}
  is simply false. Closedness under addition says that $v_1, v_2 \in V_3 \implies v_1 + v_2 \in V_3$. The attempt uses the \emph{converse} of that implication, which no subspace satisfies: a sum can land inside a subspace while neither summand does.

  A two-line counterexample in $\real^2$, with $V_3$ the $x$-axis:
  \begin{equation}
    v_1 = (1,1) \notin V_3, \quad v_2 = (0,-1) \notin V_3, \quad \text{yet} \quad v_1 + v_2 = (1,0) \in V_3.
  \end{equation}

  There is a second, subtler flaw hiding underneath. Even if one somehow knew $v_1 \in V_3$, the decomposition $v = v_1 + v_2$ is \emph{not unique} in general, so \qt{the} $v_1$ and $v_2$ are not well-defined objects to reason about. The claim would need \emph{some} decomposition with both pieces in $V_3$, not a statement about an arbitrarily chosen one.

  And the conclusion really is false, not merely unproved: the three lines of Exercise 19 give $(V_1+V_2) \cap V_3 = V_3 \neq \{0\} = (V_1 \cap V_3) + (V_2 \cap V_3)$. The identity is rescued only by a containment hypothesis such as $V_1 \subseteq V_3$, which is exactly what \eqref{eq: lemma} assumes --- and there the containment is used precisely to conclude $w = v - u \in V_3$, the step this attempt tries to get for free.
\end{ainote}
